\section{Conclusão}

Este estudo se propôs a desenvolver um modelo preditivo para o comportamento de compra de clientes em e-commerce, com foco em modelar o tempo até a próxima compra utilizando técnicas de análise de sobrevivência. Ao final da investigação, a conclusão mais significativa é a notável homogeneidade no comportamento de recompra dos clientes presentes no dataset.

A Análise Exploratória de Dados (AED) inicial já indicava que variáveis estáticas e demográficas, como gênero e idade, bem como características da transação, como categoria de produto e método de pagamento, não possuíam uma correlação clara ou forte com o abandono (\textit{churn}) dos clientes. Essa ausência de preditores simples foi confirmada e aprofundada pela modelagem de sobrevivência.

A análise não-paramétrica, através do teste Log-Rank, demonstrou que não há diferença estatisticamente significativa nas curvas de sobrevivência ao segmentar os clientes por gênero, volume de compra ou categoria de produto. Diante dessa homogeneidade, a curva de Kaplan-Meier geral (Figura \ref{fig:km_geral}) emerge como a ferramenta mais valiosa e prática deste estudo. Ela fornece um modelo robusto do comportamento geral da base de clientes, estabelecendo uma mediana de tempo para a recompra de \textbf{187 dias}. Este valor serve como um importante benchmark para a empresa, definindo uma "janela de oportunidade" para ações de reengajamento.

A modelagem semi-paramétrica com o modelo de Riscos Proporcionais de Cox corroborou estes achados. O Índice de Concordância (C-index) de 0.50 indicou que o modelo, com as variáveis disponíveis, não possui poder preditivo superior ao acaso. A análise dos coeficientes reforçou que nenhuma variável teve um impacto estatisticamente significativo no tempo de recompra, com a exceção de um efeito marginal e de pequena magnitude da variável \texttt{Returns}.

Respondendo aos objetivos do trabalho, o estudo foi bem-sucedido em modelar o comportamento de compra, e o principal fator identificado foi, paradoxalmente, a ausência de influência das variáveis testadas. A implicação prática deste resultado é que estratégias de marketing e retenção excessivamente segmentadas com base nestes critérios podem não ser eficientes. Em vez disso, uma abordagem mais ampla, baseada no tempo desde a última compra — por exemplo, campanhas direcionadas a todos os clientes que se aproximam da marca de 187 dias — pode gerar um retorno sobre o investimento mais eficaz.

Por fim, é importante reconhecer as limitações do estudo, principalmente o uso de um conjunto de dados sintético, que pode não refletir todas as nuances do comportamento real do consumidor. Trabalhos futuros poderiam se beneficiar da aplicação desta metodologia em dados reais e da inclusão de variáveis mais dinâmicas e comportamentais, como frequência de visitas ao site, tempo de navegação ou interações com campanhas de marketing, que poderiam revelar padrões não capturados pelas variáveis estáticas aqui analisadas.